% ==============================================================================
% CẤU HÌNH HỆ THỐNG & PREAMBLE CHUẨN ĐỒ ÁN
% Tương thích hoàn hảo với cả XeLaTeX (khuyên dùng trên macOS) và pdfLaTeX
% ==============================================================================
\usepackage{iftex}

\ifXeTeX
  \usepackage{fontspec}
  \setmainfont{Times New Roman}
\else
  \usepackage[utf8]{inputenc}
  \usepackage[T5,T1]{fontenc}
\fi

% ==============================================================================
% 1. VIỆT HÓA CÁC TIÊU ĐỀ MẶC ĐỊNH
% ==============================================================================
\newcommand{\chaptertitlename}{Chương}
\renewcommand{\contentsname}{Mục lục}
\renewcommand{\listfigurename}{Danh mục hình vẽ}
\renewcommand{\listtablename}{Danh mục bảng biểu}
\renewcommand{\bibname}{Tài liệu tham khảo}
\renewcommand{\figurename}{Hình}
\renewcommand{\tablename}{Bảng}

% ==============================================================================
% 2. CĂN LỀ CHUẨN ĐỒ ÁN / BÁO CÁO KHOA HỌC
%    Top: 2.5cm | Bottom: 2.5cm | Left: 3.5cm (đóng gáy) | Right: 2.0cm
% ==============================================================================
\usepackage{geometry}
\geometry{
    a4paper,
    top=25mm,
    bottom=25mm,
    left=35mm,
    right=20mm,
    includehead,
    includefoot
}

% Giãn dòng chuẩn 1.3 & thụt lề
\usepackage{setspace}
\setstretch{1.3}
\setlength{\parindent}{1cm}
\setlength{\parskip}{6pt}

% ==============================================================================
% 3. GÓI TOÁN HỌC & HÌNH ẢNH
% ==============================================================================
\usepackage{amsmath,amssymb,amsthm}
\usepackage{graphicx}
\graphicspath{{figures/}} % Đường dẫn đến thư mục hình ảnh
\usepackage{xcolor}
\usepackage{tikz}
\usetikzlibrary{arrows.meta, positioning, calc}
\usepackage{booktabs}
\usepackage{array}

% ==============================================================================
% 4. ĐỊNH DẠNG TIÊU ĐỀ CHƯƠNG CHUẨN (KHÔNG CẦN CÀI THÊM GÓI PHỤ)
% ==============================================================================
\makeatletter
\def\@makechapterhead#1{%
  \vspace*{5\p@}%
  {\parindent \z@ \centering \normalfont
    \ifnum \c@secnumdepth >\m@ne
      \large\bfseries \MakeUppercase{\chaptertitlename}\ \thechapter
      \par\nobreak
      \vskip 8\p@
    \fi
    \interlinepenalty\@M
    \LARGE \bfseries \MakeUppercase{#1}\par\nobreak
    \vskip 22\p@
  }}
\def\@makeschapterhead#1{%
  \vspace*{5\p@}%
  {\parindent \z@ \centering \normalfont
    \interlinepenalty\@M
    \LARGE \bfseries \MakeUppercase{#1}\par\nobreak
    \vskip 22\p@
  }}
\makeatother

% ==============================================================================
% 5. ĐỊNH DẠNG MÔI TRƯỜNG ĐỊNH LÝ / BỔ ĐỀ / BÀI TẬP
% ==============================================================================
\theoremstyle{plain}
\newtheorem{theorem}{Định lý}[chapter]
\newtheorem{lemma}[theorem]{Bổ đề}
\newtheorem{corollary}[theorem]{Hệ quả}
\newtheorem{proposition}[theorem]{Mệnh đề}

\theoremstyle{definition}
\newtheorem{definition}[theorem]{Định nghĩa}
\newtheorem{example}[theorem]{Ví dụ}
\newtheorem{exercise}{Bài tập}[chapter]

\theoremstyle{remark}
\newtheorem{remark}[theorem]{Nhận xét}

\renewcommand{\qedsymbol}{$\blacksquare$}

% ==============================================================================
% 6. HEADER & FOOTER (FANCYHDR)
% ==============================================================================
\usepackage{fancyhdr}
\pagestyle{fancy}
\fancyhf{}
\setlength{\headheight}{15pt}
\fancyhead[L]{\small\slshape Đề tài 5: Khôi phục tín hiệu thưa}
\fancyhead[R]{\small\slshape\nouppercase{\leftmark}}
\fancyfoot[C]{\thepage}
\renewcommand{\headrulewidth}{0.5pt}
\renewcommand{\footrulewidth}{0.3pt}

% Trang đầu chương không hiện header
\fancypagestyle{plain}{
    \fancyhf{}
    \fancyfoot[C]{\thepage}
    \renewcommand{\headrulewidth}{0pt}
    \renewcommand{\footrulewidth}{0pt}
}

% ==============================================================================
% 7. HYPERLINK & BOOKMARK
% ==============================================================================
\usepackage[
    colorlinks=true,
    linkcolor=blue!75!black,
    citecolor=teal!80!black,
    urlcolor=blue!60!black,
    bookmarksnumbered=true,
    pdfstartview=FitH
]{hyperref}
\usepackage{bookmark}
